\documentclass[11pt,reqno]{amsart}

\usepackage[T1]{fontenc}
\usepackage{lmodern}
\usepackage{microtype}
\usepackage{mathtools}
\usepackage{amssymb}
\usepackage{enumitem}
\usepackage{xcolor}
\usepackage[colorlinks=true,
  linkcolor=blue!55!black,
  citecolor=green!40!black,
  urlcolor=blue!60!black]{hyperref}

\newcommand{\Q}{\mathbf{Q}}
\newcommand{\M}{\mathcal{M}}
\newcommand{\BM}{\mathrm{BM}}
\newcommand{\RH}{\mathrm{RH}}
\newcommand{\Mod}{\operatorname{Mod}}
\newcommand{\PMod}{\operatorname{PMod}}
\newcommand{\Gr}{\operatorname{Gr}}
\newcommand{\vcd}{\operatorname{vcd}}
\newcommand{\im}{\operatorname{im}}

\DeclareMathOperator{\codim}{codim}

\newtheorem{theorem}{Theorem}[section]
\newtheorem{proposition}[theorem]{Proposition}
\newtheorem{lemma}[theorem]{Lemma}
\newtheorem{corollary}[theorem]{Corollary}
\theoremstyle{definition}
\newtheorem{definition}[theorem]{Definition}
\theoremstyle{remark}
\newtheorem{remark}[theorem]{Remark}

\title[The lowest-weight cohomology of $\M_{6,6}$]
  {The lowest-weight cohomology of $\mathcal{M}_{6,6}$ is Tate}
\author{GPT-5.6, Claude Opus 5, Qwen 3.8, Grok 4.6, and Daisuke Ken Sakai}
\date{26 August 2026}

\begin{document}

\begin{abstract}
We determine the Hodge type of the lowest-weight part of
$H^{26}(\M_{6,6};\Q)$.  We prove that
\[
  W_{26}H^{26}(\M_{6,6};\Q)
  =\RH^{26}(\M_{6,6};\Q)
  \cong \Q(-13)^{\oplus r}
\]
for some $r\geq 0$.  Equivalently,
$W_{-16}H^{\BM}_{16}(\M_{6,6};\Q)\cong\Q(8)^{\oplus r}$.
The proof combines the description by Canning--Larson--Payne of the
primitive quotient of the lowest-weight cohomology of $\M_{g,n}$ with
Harer's computation of the virtual cohomological dimension.  The primitive
quotient vanishes in the case at hand, as does the contribution pulled back
from $\M_{6,5}$.  The remaining classes are products of cotangent-line
classes with lowest-weight classes on $\M_{6,5}$; the latter are algebraic by
the sharp genus-six Chow--K\"unneth generation result.  The argument does not
determine the multiplicity $r$.
\end{abstract}

\maketitle

\section{Introduction}

Let $\M_{g,n}$ be the Deligne--Mumford stack of smooth complex curves of
genus $g$ with $n$ ordered marked points.  We work throughout with rational
cohomology and with the mixed Hodge structures carried by the cohomology of
complex algebraic stacks.  The group considered in this paper is
\[
  W_{26}H^{26}(\M_{6,6};\Q),
\]
the lowest-weight part of the degree-$26$ cohomology.  Its Hodge type is the
open contribution associated with the pair $(g,n)=(6,6)$ in the degree-$16$
boundary induction for the moduli stack of stable curves.  Indeed,
$\dim_{\mathbf C}\M_{6,6}=21$, and Poincar\'e duality identifies it, after a
Tate twist, with the weight-$-16$ part of degree-$16$ Borel--Moore homology.

Our main result is the following.

\begin{theorem}\label{thm:main}
There is a nonnegative integer $r$ such that
\begin{align*}
  W_{26}H^{26}(\M_{6,6};\Q)
    &=\RH^{26}(\M_{6,6};\Q)
      \cong \Q(-13)^{\oplus r},\\
  W_{-16}H^{\BM}_{16}(\M_{6,6};\Q)
    &\cong \Q(8)^{\oplus r},\\
  \Gr^{W}_{16}H^{16}_c(\M_{6,6};\Q)
    &\cong \Q(-8)^{\oplus r}.
\end{align*}
Here $\RH^*$ denotes tautological cohomology.  The assertion includes the
possibility $r=0$.
\end{theorem}

The proof is short once two published structural results are put together at
their exact numerical limits.  Canning, Larson, and Payne describe the
quotient of $W_kH^k(\M_{g,n})$ by classes pulled back from one fewer marking
and by their products with cotangent-line classes
\cite[Lemma~3.1]{CanningLarsonPayne}.  At $(g,n,k)=(6,6,26)$, the quotient is
built from degree-$20$ cohomology of symplectic local systems on $\M_6$.
Harer's formula gives $\vcd(\Mod_6)=19$, so every such group vanishes.  The
part pulled back from $\M_{6,5}$ vanishes as well, because
$\vcd(\PMod_{6,5})=25$.

It remains to determine the Hodge type of classes obtained by multiplying
pullbacks from $W_{24}H^{24}(\M_{6,5})$ by cotangent-line classes.  This is
where the sharp endpoint of the known Chow--K\"unneth generation range is
essential: the genus-six bound is five marked points
\cite[Table~1]{CanningLarsonPayne}.  Proposition~4.5 of
\cite{CanningLarsonPayne} therefore shows that
$W_{24}H^{24}(\M_{6,5})$ is tautological and hence of type $(12,12)$.
Multiplication by a cotangent-line class gives type $(13,13)$.

The argument uses virtual cohomological dimension only in degrees strictly
above the relevant bounds: degree $20>19$ on $\M_6$ and degree $26>25$ on
$\M_{6,5}$.  In particular, it does not infer any vanishing from the equality
$26=\vcd(\PMod_{6,6})$.  We also stress that Theorem~\ref{thm:main} determines
the Hodge type, but not the dimension, of the group.

\section{Conventions and the forgetful subspaces}

The Tate Hodge structure $\Q(m)$ has weight $-2m$.  For a smooth complex
algebraic stack $X$, the weights on $H^j(X;\Q)$ are at least $j$; consequently
$W_jH^j(X;\Q)$ is a pure Hodge structure of weight $j$.  We use the increasing
weight filtration.  Our conventions agree with those of
\cite{CanningLarsonPayne}.  The basic mixed Hodge theory is due to Deligne
\cite{DeligneHodgeII}; rational cohomology and Poincar\'e duality for
Deligne--Mumford stacks may be obtained from a finite level cover and are also
discussed in \cite{BehrendStacks}.

For $1\leq i\leq n$, let
\[
  \pi_i\colon \M_{g,n}\longrightarrow \M_{g,n-1}
\]
be the map forgetting the $i$th marking, and let
$\psi_i=c_1(L_i)\in H^2(\M_{g,n};\Q)$ be the first Chern class of the
corresponding cotangent line.  Following \cite[Sections~2--3]
{CanningLarsonPayne}, define
\begin{align}
  \Phi^k_{g,n}
    &:=\sum_{i=1}^n
      \pi_i^*W_kH^k(\M_{g,n-1};\Q),\label{eq:Phi}\\
  \Psi^k_{g,n}
    &:=\sum_{i=1}^n \psi_i\smile\Phi^{k-2}_{g,n}.
      \label{eq:Psi}
\end{align}
Thus $\Phi^k_{g,n}$ consists of lowest-weight classes inherited from one
fewer marking, while $\Psi^k_{g,n}$ is generated by such classes after one
multiplication by a cotangent-line class.

Every algebraic cycle of codimension $a$ has cohomology class of Hodge type
$(a,a)$ and weight $2a$.  In particular,
\[
  \RH^k(\M_{g,n};\Q)\subseteq W_kH^k(\M_{g,n};\Q),
\]
with both sides zero in odd degree for the tautological contribution.  The
classes $\psi_i$ are tautological and have type $\Q(-1)$.

\section{The primitive quotient and algebraicity at five markings}

We recall the two results about pure cohomology that will be used in the
proof.  Let $\mathbb V_\lambda$ denote the irreducible symplectic local system
on $\M_g$ indexed by a partition
$\lambda=(\lambda_1\geq\cdots\geq\lambda_g\geq0)$, and put
$|\lambda|=\sum_i\lambda_i$.  If $|\lambda|=n$, write
$V_{\lambda^{\mathsf T}}$ for the corresponding rational representation of
the symmetric group; as a Hodge structure it has type $(0,0)$.

\begin{proposition}[Canning--Larson--Payne]\label{prop:primitive}
If $g\geq2$, then there is an isomorphism of pure rational Hodge structures
\begin{equation}\label{eq:primitive}
 \frac{W_kH^k(\M_{g,n};\Q)}{\Phi^k_{g,n}+\Psi^k_{g,n}}
 \cong
 \bigoplus_{|\lambda|=n}
 W_kH^{k-n}(\M_g;\mathbb V_\lambda)
 \otimes V_{\lambda^{\mathsf T}}.
\end{equation}
\end{proposition}

\begin{proof}
This is \cite[Lemma~3.1(a)]{CanningLarsonPayne}.  We briefly recall the
geometry behind the formula.  If $\pi\colon\mathcal C\to\M_g$ is the universal
closed curve, then $\M_{g,n}$ is the complement of the pairwise diagonals in
the $n$-fold fiber power $\mathcal C^n$ over $\M_g$.  Restriction from the
fiber power surjects onto the lowest-weight cohomology of $\M_{g,n}$.  The
smooth-proper Leray decomposition for $\mathcal C^n\to\M_g$ has K\"unneth
indices in $\{0,1,2\}^n$.  The summands with a zero index are generated by
forgetful pullbacks.  More precisely, the identification of the index-two
summands modulo the zero-index summands is made by the relative
Chow--K\"unneth projector $\pi_2$ of Petersen--Tavakol--Yin
\cite[Sections~5.1 and 5.2.2]{PetersenTavakolYin}; it is not inferred merely
from the restriction of a cotangent-line class to a fiber.  The image of this
projector is generated by the cotangent-line contribution in
\eqref{eq:Psi}.  The all-index-one term is the $n$th tensor power of
$R^1\pi_*\Q$; quotienting by contractions coming from the diagonals leaves
precisely the symplectic local systems of weight $n$ in
\eqref{eq:primitive}.
\end{proof}

The second input identifies lowest-weight cohomology with algebraic classes
at the endpoint of the presently known genus-six range.

\begin{proposition}\label{prop:algebraic}
For every $k$,
\[
  W_kH^k(\M_{6,5};\Q)=\RH^k(\M_{6,5};\Q).
\]
In particular, for some $a\geq0$,
\[
  W_{24}H^{24}(\M_{6,5};\Q)
  \cong \Q(-12)^{\oplus a}.
\]
\end{proposition}

\begin{proof}
Table~1 of Canning--Larson--Payne records the inclusive genus-six bound
$c(6)=5$ for the Chow--K\"unneth generation property and tautological Chow
ring of the open stack $\M_{6,n}$.  Their Lemma~4.3 and Proposition~4.5 then
give
\cite[Table~1, Lemma~4.3, and Proposition~4.5]{CanningLarsonPayne}
\[
  W_kH^k(\M_{6,5};\Q)=\RH^k(\M_{6,5};\Q)
\]
for every $k$.  In degree $24$, every
class on the right is the class of a rational linear combination of
codimension-$12$ tautological cycles, so the Hodge structure is concentrated
in type $(12,12)$.  It is therefore a finite direct sum of $\Q(-12)$.
\end{proof}

\section{Virtual cohomological dimension}

Let $\Mod_g$ be the mapping class group of a closed oriented surface of genus
$g$, and let $\PMod_{g,n}$ be the pure mapping class group with $n$ ordered
punctures.  The analytic stacks $\M_g$ and $\M_{g,n}$ have the homotopy types
of the corresponding classifying stacks, since Teichm\"uller space is
contractible.  Hence cohomology of a rational local system agrees with group
cohomology of the associated rational module.

\begin{lemma}\label{lem:vcdcoeff}
Let $\Gamma$ be a virtually torsion-free group of virtual cohomological
dimension $d$.  If $L$ is a rational $\Gamma$-module, then
\[
  H^q(\Gamma;L)=0\qquad(q>d).
\]
\end{lemma}

\begin{proof}
Choose a torsion-free normal subgroup $\Gamma'\triangleleft\Gamma$ of finite
index with cohomological dimension $d$.  Then
$H^q(\Gamma';L)=0$ for $q>d$.  In the Hochschild--Serre spectral sequence for
$1\to\Gamma'\to\Gamma\to F\to1$, the group $F$ is finite.  Since the
coefficient field has characteristic zero, taking $F$-invariants is exact and
$H^p(F;U)=0$ for $p>0$ and every rational $F$-module $U$.  The spectral
sequence therefore gives
\[
  H^q(\Gamma;L)\cong H^q(\Gamma';L)^F,
\]
which vanishes for $q>d$.  See also \cite[Chapter~VII]{BrownGroups}.
\end{proof}

Harer computed \cite[Theorem~4.1]{HarerVCD}
\begin{equation}\label{eq:vcd}
 \vcd(\Mod_g)=4g-5,
 \qquad
 \vcd(\PMod_{g,n})=4g-4+n\quad(n>0),
\end{equation}
for $g\geq2$.  The two instances needed below are therefore
\[
  \vcd(\Mod_6)=19,
  \qquad
  \vcd(\PMod_{6,5})=25.
\]

\begin{corollary}\label{cor:vcdvanishing}
The following vanishings hold:
\begin{enumerate}[label=\textup{(\roman*)}]
\item $H^{20}(\M_6;\mathbb V)=0$ for every finite-dimensional rational local
  system $\mathbb V$;
\item $H^{26}(\M_{6,5};\Q)=0$.
\end{enumerate}
\end{corollary}

\begin{proof}
Apply Lemma~\ref{lem:vcdcoeff} and \eqref{eq:vcd}.  The first degree is
$20>19$, and the second is $26>25$.
\end{proof}

\section{Proof of the main theorem}

We first determine the entire lowest-weight group in ordinary cohomology.

\begin{proposition}\label{prop:psi}
One has
\[
  W_{26}H^{26}(\M_{6,6};\Q)=\Psi^{26}_{6,6}.
\]
\end{proposition}

\begin{proof}
Apply Proposition~\ref{prop:primitive} with $(g,n,k)=(6,6,26)$.  Its
right-hand side is
\[
  \bigoplus_{|\lambda|=6}
  W_{26}H^{20}(\M_6;\mathbb V_\lambda)
  \otimes V_{\lambda^{\mathsf T}},
\]
which vanishes by Corollary~\ref{cor:vcdvanishing}(i).  It follows that
\begin{equation}\label{eq:PhiPsi}
  W_{26}H^{26}(\M_{6,6};\Q)
  =\Phi^{26}_{6,6}+\Psi^{26}_{6,6}.
\end{equation}
On the other hand, Corollary~\ref{cor:vcdvanishing}(ii) and the definition
\eqref{eq:Phi} give
\[
  \Phi^{26}_{6,6}
  =\sum_{i=1}^6\pi_i^*W_{26}H^{26}(\M_{6,5};\Q)=0.
\]
Substituting this in \eqref{eq:PhiPsi} proves the proposition.
\end{proof}

\begin{proof}[Proof of Theorem~\ref{thm:main}]
By Proposition~\ref{prop:algebraic}, each summand used to generate
$\Phi^{24}_{6,6}$ is a quotient of a finite direct sum of copies of
$\Q(-12)$.  Hence $\Phi^{24}_{6,6}$ is concentrated in Hodge type $(12,12)$.
Cup product with $\psi_i$, an algebraic class of type $(1,1)$, is a morphism
of Hodge structures
\[
  \Q(-1)\otimes\Phi^{24}_{6,6}
  \longrightarrow W_{26}H^{26}(\M_{6,6};\Q).
\]
It follows that $\Psi^{26}_{6,6}$ is concentrated in type $(13,13)$.  By
Proposition~\ref{prop:psi}, the same is true of the entire group
$W_{26}H^{26}(\M_{6,6};\Q)$.  A rational Hodge structure concentrated in
type $(13,13)$ is, noncanonically, a direct sum of copies of $\Q(-13)$.
Thus
\[
  W_{26}H^{26}(\M_{6,6};\Q)\cong\Q(-13)^{\oplus r}
\]
for some $r\geq0$.

Moreover, Proposition~\ref{prop:psi} expresses every class in this group as
a sum of products of tautological classes.  Thus
$W_{26}H^{26}(\M_{6,6};\Q)\subseteq\RH^{26}(\M_{6,6};\Q)$.  The reverse
inclusion holds because every codimension-$13$ tautological cycle class has
weight $26$.  This proves the first line of the theorem.

It remains to translate the result into degree-$16$ homology.  The complex
dimension of $\M_{6,6}$ is
\[
  d=3g-3+n=21.
\]
Poincar\'e duality for the smooth Deligne--Mumford stack $\M_{6,6}$ gives an
isomorphism of mixed Hodge structures
\[
  H^{\BM}_{16}(\M_{6,6};\Q)
  \cong H^{2d-16}(\M_{6,6};\Q)(d)
  =H^{26}(\M_{6,6};\Q)(21).
\]
Taking the indicated lowest-weight piece yields
\[
  W_{-16}H^{\BM}_{16}(\M_{6,6};\Q)
  \cong W_{26}H^{26}(\M_{6,6};\Q)(21)
  \cong \Q(-13)^{\oplus r}(21)
  =\Q(8)^{\oplus r}.
\]
Finally, Borel--Moore homology and compactly supported cohomology are dual as
mixed Hodge structures.  With the standard dual weight filtration, this gives
\[
  \Gr^W_{16}H_c^{16}(\M_{6,6};\Q)
  \cong
  \operatorname{Hom}_{\Q}\!\left(
    \Gr^W_{-16}H^{\BM}_{16}(\M_{6,6};\Q),\Q(0)
  \right).
\]
The weights on $H^{\BM}_{16}(\M_{6,6};\Q)$ are at least $-16$, so its
weight-$-16$ graded piece is $W_{-16}H^{\BM}_{16}$.  Since
$\Q(8)^\vee=\Q(-8)$, the preceding Borel--Moore calculation therefore gives
\[
  \Gr^W_{16}H_c^{16}(\M_{6,6};\Q)
  \cong\Q(-8)^{\oplus r},
\]
as claimed.
\end{proof}

\section{Remarks on scope}

\begin{remark}[What is proved]
The proof determines the Hodge type of the full lowest-weight group, not only
of a specified tautological subspace.  In fact, it proves the stronger
generation statement
\[
  W_{26}H^{26}(\M_{6,6};\Q)=\Psi^{26}_{6,6}
  =\RH^{26}(\M_{6,6};\Q).
\]
Thus no non-Tate summand can survive in this group.
\end{remark}

\begin{remark}[What is not proved]
The multiplicity $r$ is not computed.  Nor does the argument assert an
integral result or identify the cohomology of the stack with that of its
coarse moduli space.  All statements are about rational mixed Hodge
structures on Deligne--Mumford stacks over $\mathbf C$.
\end{remark}

\begin{remark}[Why the numerical endpoint matters]
The only input that determines Hodge type, rather than merely generation, is
Proposition~\ref{prop:algebraic} at $n=5$.  The bound $n\leq5$ in genus six is
inclusive and sharp for the present argument.  No Chow--K\"unneth generation
statement is applied directly to $\M_{6,6}$.
\end{remark}

\section*{Acknowledgments}

The human author credits the principal proof development, mathematical
analysis, drafting, and review underlying this manuscript to the artificial
intelligence collaborators named in the byline.

\begin{thebibliography}{99}

\bibitem{BehrendStacks}
K.~Behrend,
\emph{Cohomology of stacks},
in \emph{Intersection theory and moduli}, ICTP Lecture Notes, vol.~XIX,
Abdus Salam International Centre for Theoretical Physics, Trieste, 2004,
pp.~249--294.

\bibitem{BrownGroups}
K.~S. Brown,
\emph{Cohomology of groups},
Graduate Texts in Mathematics, vol.~87, Springer-Verlag, New York, 1982.

\bibitem{CanningLarsonPayne}
S.~Canning, H.~Larson, and S.~Payne,
\emph{Extensions of tautological rings and motivic structures in the
cohomology of $\overline{\mathcal M}_{g,n}$},
Forum Math. Pi \textbf{12} (2024), article no.~e23, 31 pp.,
\href{https://doi.org/10.1017/fmp.2024.24}{doi:10.1017/fmp.2024.24}.

\bibitem{DeligneHodgeII}
P.~Deligne,
\emph{Th\'eorie de Hodge. II},
Publ. Math. Inst. Hautes \'Etudes Sci. \textbf{40} (1971), 5--57,
\href{https://doi.org/10.1007/BF02684692}{doi:10.1007/BF02684692}.

\bibitem{HarerVCD}
J.~L. Harer,
\emph{The virtual cohomological dimension of the mapping class group of an
orientable surface},
Invent. Math. \textbf{84} (1986), no.~1, 157--176,
\href{https://doi.org/10.1007/BF01388737}{doi:10.1007/BF01388737}.

\bibitem{PetersenTavakolYin}
D.~Petersen, M.~Tavakol, and Q.~Yin,
\emph{Tautological classes with twisted coefficients},
Ann. Sci. \'Ec. Norm. Sup\'er. (4) \textbf{54} (2021), no.~5, 1179--1236,
\href{https://doi.org/10.24033/asens.2479}{doi:10.24033/asens.2479}.

\end{thebibliography}

\end{document}
